\section{9차시 --- 발사대 제작과 1분기 회고}
\label{sec:ch09}

본 차시는 코딩 없이 "손으로 만드는" 시간이다. 다음 차시부터 시작될
로켓 발사 시스템의 무대인 발사대 건물을 학습자가 직접 조립·점검하고,
1분기에서 배운 부품과 개념을 총정리한다. 4개 미션 모두 이론/제작
중심이며 본격 코딩은 다음 차시로 이어진다.

\subsection{미션 1 --- 발사대 건물 조립 시작!}
\label{subsec:ch09-m1}
\missionmeta{(제작 활동)}{theory}

\begin{storybox}
\sayares{알비! 드디어 발사대를 만들 날이야! 이 발사대에 LED와 부저를 달면 진짜 로켓 발사 시스템이 완성돼!}
\sayalbi{삐릿♪ 설계도대로 꼼꼼하게 만들어봐요! 다음 수업에서 멋진 발사 시스템을 달아드릴게요!}
\end{storybox}

\subsubsection*{학습 목표}
\begin{itemize}[leftmargin=1.4em]
  \item 발사대 건물 설계도대로 제작
  \item 부품 목록 확인 및 준비
  \item 뼈대 구조 먼저 조립하기
\end{itemize}

\begin{labbox}{샘플 코드 --- 9차시: 발사대 건물 제작}
\begin{minted}{python}
# 9차시: 발사대 건물 제작
# 이번 시간은 제작 미션입니다.
# 코딩 없이 손으로 발사대를 만들어요!
\end{minted}
\end{labbox}

\subsection{미션 2 --- 발사대 완성 및 점검}
\label{subsec:ch09-m2}
\missionmeta{(제작/점검)}{theory}

\begin{storybox}
\sayares{발사대가 거의 완성됐어! 단단하게 고정됐는지 확인해야 해.}
\sayalbi{삐릿- 구조물이 튼튼한지 꼭 확인하세요! 다음 시간에 여기에 LED 5개를 달 거예요!}
\end{storybox}

\subsubsection*{학습 목표}
\begin{itemize}[leftmargin=1.4em]
  \item 조립된 발사대 구조 점검
  \item LED 5개를 달 위치 미리 확인
  \item 발사대 완성 후 사진 찍기
\end{itemize}

\begin{labbox}{샘플 코드 --- 발사대 완성!}
\begin{minted}{python}
# 발사대 완성!
# 다음 시간에 LED를 달고 코딩을 시작해요.
print("발사대 건물 완성!")
\end{minted}
\end{labbox}

\subsection{미션 3 --- 발사 시스템 계획 세우기}
\label{subsec:ch09-m3}
\missionmeta{(설계 활동)}{theory}

\begin{storybox}
\sayalbi{삐릿! 다음 시간에 만들 발사 시스템을 미리 계획해볼까요?}
\end{storybox}

\subsubsection*{학습 목표}
\begin{itemize}[leftmargin=1.4em]
  \item LED 5개 순차 점등 계획
  \item 카운트다운 흐름 설계
  \item LED + 부저 + 모터 통합 시나리오 작성
\end{itemize}

\begin{labbox}{샘플 코드 --- 10$\sim$12차시 계획}
\begin{minted}{python}
# 10~12차시 계획
# LED 5개: 순차 점등 → 카운트다운
# 부저: 사이렌 + 멜로디
# DC모터: 발사 장치 작동
\end{minted}
\end{labbox}

\subsection{미션 4 --- 1$\sim$2분기 연결 이해하기}
\label{subsec:ch09-m4}
\missionmeta{(이론/회고)}{theory}

\begin{storybox}
\sayares{알비! 1분기에서 배운 것들을 정리해볼까?}
\sayalbi{삐릿♪ LED, 부저, DC모터를 모두 배웠어요! 이것들을 합치면 진짜 로켓 발사 시스템이 완성돼요!}
\end{storybox}

\subsubsection*{학습 목표}
\begin{itemize}[leftmargin=1.4em]
  \item 1분기 부품 총복습: LED, 부저, DC모터
  \item 순차·반복·시간지연·디지털출력·PWM 정리
  \item 2분기 예고: 서보모터, 초음파 센서, LCD
\end{itemize}

\begin{labbox}{샘플 코드 --- 1분기 총정리}
\begin{minted}{python}
# 1분기 총정리
print("배운 부품: LED, 부저, DC모터")
print("배운 개념: 순차, 반복, 시간지연, PWM")
\end{minted}
\end{labbox}

\begin{notebox}{교사용 메모}
9차시는 인지 부하가 낮은 "숨 고르기" 차시이다. 학습자가 손으로
제작 활동을 하면서 1분기 학습 내용을 \emph{체화}하고, 10$\sim$12차시의
통합 프로젝트(로켓 발사 시스템)에 대한 동기를 부여하는 데 초점을 둔다.
\end{notebox}
